\documentclass[11pt]{article}
\usepackage[a4paper,margin=25mm]{geometry}
\usepackage[no-math]{fontspec}
\setmainfont{Latin Modern Roman}
\usepackage{amsmath,amssymb,amsthm,xfrac,xspace,hyperref,graphicx,comment}
\excludecomment{DefTralics}
\providecommand{\C}{}
\newcommand{\ie}{i.e.,\xspace}
\newcommand{\eg}{e.g.,\xspace}
\newcommand{\etc}{etc.\xspace}
\newcommand{\cf}{cf.\xspace}
\newcommand{\resp}{resp.\xspace}
\newcommand{\smaller}{\small}
\newcommand{\Subsection}{\subsection}
\newcommand{\Subsubsection}{\subsubsection}
\newcommand{\skpt}{\smallskip}
\emergencystretch=3em
\input{author-preamble.tex}
\begin{document}
\begin{center}{\Large Stable item 1925}\end{center}
\paragraph{Source and attribution.}
Matthew Emerton, Vytautas Paškūnas, \emph{On the density of supercuspidal points of fixed regular weight in local deformation rings and global Hecke algebras}, Journal de l’École polytechnique — Mathématiques 7 (2020), publisher version, DOI 10.5802/jep.119. \href{https://jep.centre-mersenne.org/articles/10.5802/jep.119/}{Publisher article}. Authors' text: \href{https://creativecommons.org/licenses/by/4.0/}{CC BY 4.0}.
\paragraph{Editorial problem boundary.}
Prove only the two-sided growth inequality of Proposition 2.18 below for a nonzero finitely generated \(k[[K]]\)-module, where \(k\) is a finite field of characteristic \(p\), \(K\) is a uniformly powerful pro-\(p\) group of positive \(p\)-adic dimension, and \(K_n=K^{p^n}\). The canonical delta dimension, exact filtration, graded coinvariant maps, and complete Hilbert--Samuel/Hilbert--Kunz sandwich proof are retained, including the authors\' alternative Hilbert--Samuel upper bound. No arbitrary filtration or arithmetic-density theorem is selected; remaining article statements are context only.
\paragraph{Difficulty (editorial): H2.}
The noncommutative Iwasawa canonical dimension is identified with the commutative associated-graded Krull dimension. A Frobenius-ideal filtration sandwiches genuine group coinvariants between Hilbert-Samuel and Hilbert-Kunz lengths, giving a sharp exponential exponent for the specified uniform power tower.
\paragraph{Editorial transcription note.}
Original author language, mathematical content and ordering are unchanged. Publisher front matter is replaced by this independent layout wrapper. Original native body and macros remain editable; bibliography is recovered from the matching cached publisher HTML. Adjacent claims are author context, not additional corpus items. Printed mathematical slips are retained without assistant repair.
\clearpage
\input{author-body.tex}
\begin{thebibliography}{99}
\bibitem{ardakov_brown} K. Ardakov \& K. A. Brown - “Ring-theoretic properties of Iwasawa algebras: a survey” , Doc. Math. (2006), p. 7-33, Extra Vol.
\bibitem{BLGGT} T. Barnet-Lamb, T. Gee, D. Geraghty \& R. Taylor - “Potential automorphy and change of weight” , Ann. of Math. (2) 179 (2014) no. 2, p. 501-609
\bibitem{BHS} C. Breuil, E. Hellmann \& B. Schraen - “Une interprétation modulaire de la variété trianguline” , Math. Ann. 367 (2017) no. 3-4, p. 1587-1645
\bibitem{bm1} C. Breuil \& A. Mézard - “Multiplicités modulaires et représentations de GL 2 ( Z p ) et de Gal ( Q ¯ p / Q p ) en \(\ell\) = p ” , Duke Math. J. 115 (2002) no. 2, p. 205-310, With an appendix by Guy Henniart
\bibitem{BF} C. J. Bushnell \& A. Fröhlich - Gauss sums and p -adic division algebras , Lect. Notes in Math. , vol. 987 , Springer-Verlag, Berlin-New York, 1983
\bibitem{BH} C. J. Bushnell \& G. Henniart - “The essentially tame local Langlands correspondence. I” , J. Amer. Math. Soc. 18 (2005) no. 3, p. 685-710
\bibitem{BK} C. J. Bushnell \& P. C. Kutzko - The admissible dual of GL ( N ) via compact open subgroups , Annals of Math. Studies , vol. 129 , Princeton University Press, Princeton, NJ, 1993
\bibitem{BK_cover} C. J. Bushnell \& P. C. Kutzko - “Semisimple types in GL n ” , Compositio Math. 119 (1999) no. 1, p. 53-97
\bibitem{calegari-emerton} F. Calegari \& M. Emerton - “Completed cohomology—a survey” , in Non-abelian fundamental groups and Iwasawa theory , London Math. Soc. Lecture Note Ser. , vol. 393 , Cambridge Univ. Press, Cambridge, 2012, p. 239-257
\bibitem{six-authors} A. Caraiani, M. Emerton, T. Gee, D. Geraghty, V. Paškūnas \& S. W. Shin - “Patching and the p -adic local Langlands correspondence” , Camb. J. Math. 4 (2016) no. 2, p. 197-287
\bibitem{Chenevier} G. Chenevier - “Sur la densité des représentations cristallines de Gal ( \(\mathbb Q\) ¯ p / \(\mathbb Q\) p ) ” , Math. Ann. 355 (2013) no. 4, p. 1469-1525
\bibitem{CDP} P. Colmez, G. Dospinescu \& V. Paškūnas - “The p -adic local Langlands correspondence for GL 2 ( \(\mathbb Q\) p ) ” , Camb. J. Math. 2 (2014) no. 1, p. 1-47
\bibitem{analytic} J. D. Dixon, M. P. F. du Sautoy, A. Mann \& D. Segal - Analytic pro- p groups , Cambridge Studies in Advanced Math. , vol. 61 , Cambridge Univ. Press, Cambridge, 1999
\bibitem{Jacquet-one} M. Emerton - “Jacquet modules of locally analytic representations of p -adic reductive groups. I. Construction and first properties” , Ann. Sci. École Norm. Sup. (4) 39 (2006) no. 5, p. 775-839
\bibitem{em0} M. Emerton - “On the interpolation of systems of eigenvalues attached to automorphic Hecke eigenforms” , Invent. Math. 164 (2006) no. 1, p. 1-84
\bibitem{emfm} M. Emerton - “Local-global compatibility in the p -adic Langlands programme for GL 2 / \(\mathbb Q\) ” (2011), Preprint, http://math.uchicago.edu/\textasciitilde{}emerton/preprints.html
\bibitem{matt_icm} M. Emerton - “Completed cohomology and the p -adic Langlands program” , in Proceedings of the International Congress of Mathematicians (Seoul, 2014). Vol. II , Kyung Moon Sa, Seoul, 2014, p. 319-342
\bibitem{BMcycles} M. Emerton \& T. Gee - “A geometric perspective on the Breuil-Mézard conjecture” , J. Inst. Math. Jussieu 13 (2014) no. 1, p. 183-223
\bibitem{gee-newton} T. Gee \& J. Newton - “Patching and the completed homology of locally symmetric spaces” , 2017
\bibitem{gross} B. H. Gross - “Algebraic modular forms” , Israel J. Math. 113 (1999), p. 61-93
\bibitem{harris} M. Harris - “ p -adic representations arising from descent on abelian varieties” , Compositio Math. 39 (1979) no. 2, p. 177-245
\bibitem{harris_err} M. Harris - “Correction to: “ p -adic representations arising from descent on abelian varieties”” , Compositio Math. 121 (2000) no. 1, p. 105-108
\bibitem{HMS} E. Hellmann, C. M. Margerin \& B. Schraen - “Density of automorphic points” , 2018
\bibitem{HS} E. Hellmann \& B. Schraen - “Density of potentially crystalline representations of fixed weight” , Compositio Math. 152 (2016) no. 8, p. 1609-1647
\bibitem{KM} P. Kutzko \& D. Manderscheid - “On intertwining operators for GL N ( F ) , F a non-Archimedean local field” , Duke Math. J. 57 (1988) no. 1, p. 275-293
\bibitem{matsumura} H. Matsumura - Commutative ring theory , Cambridge Studies in Advanced Math. , vol. 8 , Cambridge Univ. Press, Cambridge, 1989
\bibitem{monsky} P. Monsky - “The Hilbert-Kunz function” , Math. Ann. 263 (1983) no. 1, p. 43-49
\bibitem{Nakamura2} K. Nakamura - “Deformations of trianguline B -pairs and Zariski density of two dimensional crystalline representations” , J. Math. Sci. Univ. Tokyo 20 (2013) no. 4, p. 461-568
\bibitem{Nakamura} K. Nakamura - “Zariski density of crystalline representations for any p -adic field” , J. Math. Sci. Univ. Tokyo 21 (2014) no. 1, p. 79-127
\bibitem{nsw2} J. Neukirch, A. Schmidt \& K. Wingberg - Cohomology of number fields , Grundlehren Math. Wiss. , vol. 323 , Springer-Verlag, Berlin, 2013, corrected second printing
\bibitem{thesis} V. Paškūnas - “Unicity of types for supercuspidal representations of GL N ” , Proc. London Math. Soc. (3) 91 (2005) no. 3, p. 623-654
\bibitem{cmf} V. Paškūnas - “The image of Colmez’s Montreal functor” , Publ. Math. Inst. Hautes Études Sci. 118 (2013), p. 1-191
\bibitem{duke} V. Paškūnas - “On the Breuil-Mézard conjecture” , Duke Math. J. 164 (2015) no. 2, p. 297-359
\bibitem{padicJL} V. Paškūnas - “On some consequences of a theorem of J. Ludwig” , 2018
\bibitem{iw} P. Schneider \& J. Teitelbaum - “Banach space representations and Iwasawa theory” , Israel J. Math. 127 (2002), p. 359-380
\bibitem{scholze} P. Scholze - “On the p -adic cohomology of the Lubin-Tate tower” , Ann. Sci. École Norm. Sup. (4) 51 (2018) no. 4, p. 811-863, With an appendix by Michael Rapoport
\bibitem{shen-ning} S.-N. Tung - “On the automorphy of 2 -dimensional potentially semi-stable deformation rings of G \(\mathbb Q\) p ” , 2018
\bibitem{shen-ning2} S.-N. Tung - “On the modularity of 2 -adic potentially semi-stable deformation rings” , 2019
\bibitem{venjakob} O. Venjakob - “On the structure theory of the Iwasawa algebra of a p -adic Lie group” , J. Eur. Math. Soc. (JEMS) 4 (2002) no. 3, p. 271-311
\end{thebibliography}
\end{document}
